\documentclass[11pt, twoside]{article}

\usepackage[utf8]{inputenc}
\usepackage[T1]{fontenc}

\usepackage{float}

\usepackage[letterpaper, margin=.8in, includehead, includefoot]{geometry}
\usepackage{fancyhdr}
\usepackage{lastpage}
\usepackage{graphicx}
\usepackage{float}
\usepackage{mwe}
\graphicspath{{img}}

\usepackage{amsmath}
\usepackage{amssymb}
\usepackage{mathtools}

\usepackage{tikz}
\usepackage{pgfplots}
\pgfplotsset{compat=1.16}
\usepackage{circuitikz}
\usepackage{multicol}

\setlength{\parindent}{0pt}
\setlength{\headheight}{15pt}

\def\Titulo{Problemario 2ndo Parcial}
\def\Materia{Señales y sistemas de control clásico}
\def\Profesor{Jesica Azusena Escobar Medina}
\def\Grupo{6CV6}

\begin{document}
\pagestyle{fancy}
\fancyhf{}
\fancyhead[LO,RE]{\textbf{\Titulo \space| Carlos Alberto Herrera Rangel | \Grupo}}
\fancyfoot[LE, RO]{\thepage\ of \pageref{LastPage}}
\section{Calcule las constantes $K_{p}$, $K_{v}$, $K_{a}$ y el error en estado estacionario para:}
\subsection{Inciso A}
\begin{equation*}
    \mathbf{G(s) = \frac{1}{s^{2}\left(s+12\right)};\space H(s) = Ks}
\end{equation*}
\textbf{Solución:}\\\\
Obteniendo $G(s)H(s)$:
\begin{align*}
    G(s)H(s) & = \frac{Ks}{s^{2}\left(s+12\right)} \\
    G(s)H(s) & = \frac{K}{s\left(s+12\right)}      \\
\end{align*}
Obteniendo $K_{p}$:
\begin{align*}
    K_{p} & = \lim_{s\to 0}{G(s)H(s)}                              \\
          & = \left|\frac{K}{s\left(s+12\right)}\right|_{s=0}      \\
          & = \frac{K}{(0)((0)+12)} = \frac{K}{0} = \boxed{\infty} \\
\end{align*}
Obteniendo $K_{v}$:
\begin{align*}
    K_{v} & = \lim_{s\to 0}{sG(s)H(s)}                \\
          & = \left|\frac{K}{s+12}\right|_{s=0}       \\
          & = \frac{K}{(0)+12} = \boxed{\frac{K}{12}}
\end{align*}
Obteniendo $K_{a}$:
\begin{align*}
    K_{a} & = \lim_{s\to0}{s^{2}G(s)H(s)}                    \\
          & = \left|\frac{Ks}{s+12}\right|_{s=0}             \\
          & = \frac{K(0)}{(0)+12} = \frac{0}{12} = \boxed{0}
\end{align*}
\newpage
Calculando el error en estado estacionario:
\begin{align*}
    e_{ss} & = \frac{1}{1+K_{p}} = \frac{1}{1+\infty} = \boxed{0}              \\
    e_{ss} & = \frac{1}{K_{v}} = \frac{1}{\frac{K}{12}} = \boxed{\frac{12}{K}} \\
    e_{ss} & = \frac{1}{K_{a}} = \frac{1}{0} = \boxed{\infty}
\end{align*}
\subsection{Inciso B}
\begin{equation*}
    \mathbf{G(s)=\frac{K(1+0.5s)}{s(s+1)(2s+1)(s^{2}+s+1)};\space H(s)=1}
\end{equation*}
\textbf{Solución:}\\\\
Obteniendo $G(s)H(s)$:
\begin{align*}
    G(s)H(s) & = \frac{K(1+0.5s)}{s(s+1)(2s+1)(s^{2}+s+1)}
\end{align*}
Obteniendo $K_{p}$:
\begin{align*}
    K_{p} & = \lim_{s\to 0}{G(s)H(s)}                              \\
          & = \left|\frac{K(1+0.5s)}{s(s+1)(2s+1)(s^{2}+s+1)}\right|_{s=0}      \\
          & = \frac{K(1+0.5(0))}{(0)((0)+1)(2(0)+1)((0)^{2}+(0)+1)} = \frac{K}{0} = \boxed{\infty} \\
\end{align*}
Obteniendo $K_{v}$:
\begin{align*}
    K_{v} & = \lim_{s\to 0}{sG(s)H(s)}                \\
          & = \left|\frac{K(1+0.5s)}{(s+1)(2s+1)(s^{2}+s+1)}\right|_{s=0}       \\
          & = \frac{K(1+0.5(0))}{((0+1)(2((0)+1)((0)^{2}+((0)+1))))} = \frac{0}{3} = \boxed{0}
\end{align*}
Obteniendo $K_{a}$:
\begin{align*}
    K_{a} & = \lim_{s\to0}{s^{2}G(s)H(s)}                    \\
          & = \left|\frac{K(1+0.5s)s}{(s+1)(2s+1)(s^{2}+s+1)}\right|_{s=0}             \\
          & = \frac{K(1+0.5(0))(0)}{((0)+1)((0)+1)((0)^{2}+(0)+1)} = \frac{0}{3} = \boxed{0}
\end{align*}
\newpage
Calculando el error en estado estacionario:
\begin{align*}
    e_{ss} & = \frac{1}{1+K_{p}} = \frac{1}{1+\infty} = \boxed{0}              \\
    e_{ss} & = \frac{1}{K_{v}} = \frac{1}{0} = \boxed{\infty} \\
    e_{ss} & = \frac{1}{K_{a}} = \frac{1}{0} = \boxed{\infty}
\end{align*}
\subsection{Inciso C}
\begin{equation*}
    \mathbf{G(s)=\frac{1}{s^{2}(s+12)};\space H(s)=\frac{5(s+1)}{s+5}}
\end{equation*}
\textbf{Solución:}\\\\
Obteniendo $G(s)H(s)$:
\begin{align*}
    G(s)H(s) & = \frac{5(s+1)}{(5+s)(s^{2}+s+12)} \\
\end{align*}
Obteniendo $K_{p}$:
\begin{align*}
    K_{p} & = \lim_{s\to 0}{G(s)H(s)}                              \\
          & = \left|\frac{5(s+1)}{(5+s)(s^{2}+s+12)}\right|_{s=0}      \\
          & = \frac{5((0)+1)}{(5+(0))((0)^{2}+(0)+12)} = \boxed{\frac{1}{12}} \\
\end{align*}
Obteniendo $K_{v}$:
\begin{align*}
    K_{v} & = \lim_{s\to 0}{sG(s)H(s)}                \\
          & = \left|\frac{5s(s+1)}{(5+s)(s^{2}+s+12)}\right|_{s=0}       \\
          & = \frac{5(0)((0)+1)}{(5+(0))((0)^{2}+(0)+12)} = \frac{0}{60} = \boxed{0}
\end{align*}
Obteniendo $K_{a}$:
\begin{align*}
    K_{a} & = \lim_{s\to0}{s^{2}G(s)H(s)}                    \\
          & = \left|\frac{5s^{2}(s+1)}{(5+s)(s^{2}+s+12)}\right|_{s=0}             \\
          & = \frac{5(0)^{2}((0)+1)}{(5+(0))((0)^{2}+(0)+12)} = \frac{0}{60} = \boxed{0}
\end{align*}
\newpage
Calculando el error en estado estacionario:
\begin{align*}
    e_{ss} & = \frac{1}{1+K_{p}} = \frac{1}{1+\frac{1}{12}} = \boxed{\frac{12}{13}}              \\
    e_{ss} & = \frac{1}{K_{v}} = \frac{1}{0} = \boxed{\infty} \\
    e_{ss} & = \frac{1}{K_{a}} = \frac{1}{0} = \boxed{\infty}
\end{align*}
\subsection{Inciso D}
\begin{equation*}
    \mathbf{G(s)=\frac{K(2s+1)(4s+1)}{s^{2}(s^{2}+s+1)};\space H(s)=1}
\end{equation*}
\textbf{Solución:}\\\\
Obteniendo $G(s)H(s)$:
\begin{align*}
    G(s)H(s) & = \frac{K(2s+1)(4s+1)}{s^{2}(s^{2}+s+1)} \\
\end{align*}
Obteniendo $K_{p}$:
\begin{align*}
    K_{p} & = \lim_{s\to 0}{G(s)H(s)}                              \\
          & = \left|\frac{K(2s+1)(4s+1)}{s^{2}(s^{2}+s+1)}\right|_{s=0}      \\
          & = \frac{K(2(0)+1)(4(0)+1)}{(0)^{2}((0)^{2}+(0)+1)} = \frac{K+1}{0} = \boxed{\infty} \\
\end{align*}
Obteniendo $K_{v}$:
\begin{align*}
    K_{v} & = \lim_{s\to 0}{sG(s)H(s)}                \\
          & = \left|\frac{K(2s+1)(4s+1)}{s(s^{2}+s+1)}\right|_{s=0}       \\
          & = \frac{K(2(0)+1)(4(0)+1)}{(0)((0)^{2}+(0)+1)} = \frac{K+1}{0} = \boxed{\infty}
\end{align*}
Obteniendo $K_{a}$:
\begin{align*}
    K_{a} & = \lim_{s\to0}{s^{2}G(s)H(s)}                    \\
          & = \left|\frac{K(2s+1)(4s+1)}{s^{2}+s+1}\right|_{s=0}             \\
          & = \frac{K(2(0)+1)(4(0)+1)}{(0)^{2}+(0)+1} = \frac{K+1}{1} = \boxed{K+1}
\end{align*}
\newpage
Calculando el error en estado estacionario:
\begin{align*}
    e_{ss} & = \frac{1}{1+K_{p}} = \frac{1}{1+\infty} = \boxed{0}              \\
    e_{ss} & = \frac{1}{K_{v}} = \frac{1}{\infty} = \boxed{0} \\
    e_{ss} & = \frac{1}{K_{a}} = \frac{1}{K+1} = \boxed{\frac{1}{K+1}}
\end{align*}
\subsection{Inciso E}
\begin{equation*}
    \mathbf{G(s)=\frac{1000}{s(s+10)(s+100)};\space H(s)=1}
\end{equation*}
\textbf{Solución:}\\\\
Obteniendo $G(s)H(s)$:
\begin{align*}
    G(s)H(s) & = \frac{1000}{s(s+10)(s+100)} \\
\end{align*}
Obteniendo $K_{p}$:
\begin{align*}
    K_{p} & = \lim_{s\to 0}{G(s)H(s)}                              \\
          & = \left|\frac{1000}{s(s+10)(s+100)}\right|_{s=0}      \\
          & = \frac{1000}{(0)((0)+10)((0)+100)} = \frac{1000}{0} = \boxed{\infty} \\
\end{align*}
Obteniendo $K_{v}$:
\begin{align*}
    K_{v} & = \lim_{s\to 0}{sG(s)H(s)}                \\
          & = \left|\frac{1000}{(s+10)(s+100)}\right|_{s=0}       \\
          & = \frac{1000}{((0)+10)((0)+100)} = \frac{1000}{1000} = \boxed{1}
\end{align*}
Obteniendo $K_{a}$:
\begin{align*}
    K_{a} & = \lim_{s\to0}{s^{2}G(s)H(s)}                    \\
          & = \left|\frac{1000s}{(s+10)(s+100)}\right|_{s=0}             \\
          & = \frac{1000(0)}{((0)+10)((0)+100)} = \frac{0}{0} = \boxed{0}
\end{align*}
\newpage
Calculando el error en estado estacionario:
\begin{align*}
    e_{ss} & = \frac{1}{1+K_{p}} = \frac{1}{1+\infty} = \boxed{0}              \\
    e_{ss} & = \frac{1}{K_{v}} = \frac{1}{1} = \boxed{1} \\
    e_{ss} & = \frac{1}{K_{a}} = \frac{1}{0} = \boxed{\infty}
\end{align*}
\section{¿Qué valor de $\mathbf{K_{v}}$ se requiere para que $\mathbf{e_{ss}\le 30\%}$?}
\begin{equation*}
    \mathbf{G(s) = \frac{1}{s(s+50)}; \space H(s)=K}
\end{equation*}
\textbf{Solución:}\\\\
Obteniendo $G(s)H(s)$:
\begin{align*}
    G(s)H(s) &= \frac{1}{s(s+10)}\cdot K\\
    &= \boxed{\frac{K}{s(s+10)}}
\end{align*}
A partir de $e_{ss}$:
\begin{align*}
    e_{ss} &= \frac{1}{K_{v}} \Rightarrow K_{v} = \frac{1}{e_{ss}}\\
    K_{v} &= \frac{1}{0.3\bar{3}} = 3.\bar{03}
\end{align*}
\section{Determine si los siguientes sistemas son estables:}
\subsection{Inciso A}
\begin{equation*}
    \mathbf{s^{7} + 3s^{6} + 3s^{5} + s^{4} + s^{3} + 3s^{2} + 3s + 10}
\end{equation*}
\textbf{Solucion:}\\\\
\begin{equation*}
    \begin{array}{c|c c c c }
        s^{7} & 1 & 3 & 1 & 3\\
        s^{6} & 3 & 1 & 1 & 10\\
        s^{5} & \\
        s^{4} & \\
        s^{3} & \\
        s^{2} & \\
        s^{1} & \\
        s^{0} & \\
    \end{array}
\end{equation*}
\subsection{Inciso B}
\begin{equation*}
    \mathbf{s^{4} + s^{3} + s^{2} + s + 3}
\end{equation*}
\subsection{Inciso C}
\begin{equation*}
    \mathbf{2s^{4} + s^{3} + 3s^{2} + 5s + 10}
\end{equation*}
\section{Determine el valor de $\mathbf{K}$ para que los siguientes sistemas sean estables:}
\subsection{Inciso A}
\begin{equation*}
    \mathbf{s^{4} + 3s^{2} + Ks^{2} + 2s + 1 = 0}
\end{equation*}
\subsection{Inciso B}
\begin{equation*}
    \mathbf{s^{3} + 3408.3s^{2} + 1204000s + 1.5 \times 10^{-4}K = 0}
\end{equation*}
\subsection{Inciso C}
\begin{equation*}
    \mathbf{s^{3} + 3Ks^{2} + (K + 2)s + 4 = 0}
\end{equation*}
\end{document}
